% ============================================================================
% sec-classification.tex — Fan-type theorems; the placeability classification
% programme on K_{2k+1}; criterion-level results.
% ===========================================================================
\section{Fan-type theorems, the placeability classification programme, and
criterion-level results}\label{sec:classification}

\subsection{Fan-type theorems beyond the block-level degree bound}
\label{subsec:astar}

Throughout this subsection $G$ is a connected graph, $n=|V(G)|$, and
$\dEV(v)$ denotes the E-degree of $v$. Recall that a vertex with
$\dEV\ge 4$ has even degree in $G$; call such a vertex \emph{high}, and
call a vertex with $1\le \dEV\le 3$ \emph{low}. All published results in
the Fan line \cite{Fan05,BS21,FHZ23} require the blocks of $\EV(G)$ to have
maximum degree at most $3$ (the triangle conditions of the three differ:
\cite{Fan05} asks every block of $\EV(G)$ to be triangle-free, \cite{FHZ23}
that $\EV(G)$ contain no pair of triangles, and the class $\mathcal{G}$ of
\cite{BS21} allows one triangular block per component). This block
condition is implied by, but does not imply, the absence of high vertices:
a cut vertex of $\EV(G)$ may belong to arbitrarily many blocks of maximum
degree at most $3$ (in $K_{1,5}$ with one pendant leaf attached at each of
its six vertices, $\EV(G)=K_{1,5}$ has a vertex of E-degree $5$ while every
block is a single edge); what is new below is therefore a block of
$\EV(G)$ of maximum degree at least $4$, not a high vertex as such.
Theorem~\ref{thm:claAstar} allows high vertices, under an odd-degree
compensation condition (A2) together with an alignment condition (A4) at
one high vertex; Theorem~\ref{thm:lasep} removes the alignment condition
when distinct high vertices have no common EV-neighbour, and
Remark~\ref{rem:claCO} weakens separation to a parity condition on the
profiles of common EV-neighbours; Theorem~\ref{thm:clatwohigh} treats the
case of exactly two high vertices that share EV-neighbours, and
Corollary~\ref{cor:clatwohigh} collects the at-most-two-high layer,
leaving only a corner open there, subject to five simultaneous constraints
at every common EV-neighbour of the two high vertices; the band-free
tightness theorem (Remark~\ref{rem:clagap}) empties this corner on the
constructed family proper whenever the counting gap is at most two.

\begin{theorem}[alignment at a high vertex]\label{thm:claAstar}
Let $G$ be a connected graph such that
\begin{enumerate}[label={},leftmargin=2.6em]
\item[(A1)] $\EV(G)$ is triangle-free;
\item[(A2)] every high vertex $v$ has odd $\dEV(v)$, and every EV-neighbour
$w$ of $v$ satisfies $\dEV(w)\le 3$;
\item[(A4)] there is a high vertex $u$ such that every even vertex
$x\notin\{u\}\cup N_{\EV}(u)$ satisfies
$\dEV(x)-|N_{\EV}(x)\cap N_{\EV}(u)|\le 3$, where $N_{\EV}(x)$ denotes the
set of EV-neighbours of $x$.
\end{enumerate}
Then $\pn(G)\le \floor{n/2}$.
\end{theorem}

Here $N_{\EV}(x)$ denotes the set of EV-neighbours of $x$, i.e.\ its
neighbours inside $\EV(G)$.\footnote{The numbering skips (A3): a condition
so numbered in an earlier draft of this material was absorbed into the
separation hypothesis of Theorem~\ref{thm:lasep}, and the remaining labels
were kept unchanged.}

\begin{proof}
Let $F=\{uv : v\in N_{\EV}(u)\}$ be the $E$-star at $u$, $|F|=k$, and let
$\tilde G=G-F$. We first analyse the parity and structure of $\tilde G$.

\emph{Step 1 (structure of $\tilde G$).}
Every edge of $F$ joins two even-degree vertices of $G$, so odd-degree
vertices of $G$ lose no edges and stay odd. The vertex $u$ loses $k$ edges,
an odd number, so $u$ is odd in $\tilde G$. Each $v\in N_{\EV}(u)$ loses
exactly one edge, so $v$ is odd in $\tilde G$. Every other vertex keeps its
degree. Hence
$V(\EV(\tilde G))=\{x\notin\{u\}\cup N_{\EV}(u) : \deg_G(x)\text{ even}\}$,
and since $\EV(\tilde G)$ is an induced subgraph of $\EV(G)$, it is
triangle-free by (A1).

Moreover every $x\in V(\EV(\tilde G))$ has $\dEV$ in $\tilde G$ equal to
$\dEV(x)-|N_{\EV}(x)\cap N_{\EV}(u)|\le 3$ by (A4). Thus \emph{every}
vertex of $\EV(\tilde G)$ has E-degree at most $3$, so every block of
$\EV(\tilde G)$ is triangle-free of maximum degree at most $3$. Applying
Theorem~\ref{int:thm-fan} to each component of $\tilde G$ (the hypotheses are
component-wise) and summing with Lemma~\ref{int:lem-sum},
\begin{equation}\label{eq:clatildebound}
\pn(\tilde G)\;\le\;\textstyle\sum_c \floor{n_c/2}\;\le\;\floor{n/2},
\end{equation}
where $n_c$ are the component orders. Fix a decomposition $\tilde D$ of
$\tilde G$ with $|\tilde D|\le\floor{n/2}$.

\emph{Step 2 (hosts).}
The neighbours of $u$ in $\tilde G$ are precisely its odd-degree neighbours
in $G$; they are untouched and odd in $\tilde G$, so by
Lemma~\ref{int:lem-LB1}(ii) each has $\tilde D\ge 1$.

\emph{Step 3 (reinsertion of the addible part).}
Apply Lemma~\ref{int:lem-fan7} with $B=F$, $h=k$, and specified edge $uv_1$
(any fixed element of $F$). The hypothesis of Lemma~\ref{int:lem-fan7} is
satisfied: every neighbour of $u$ has $\tilde D\ge 1$ by Step 2 (this
holds whether the hypothesis is read over the neighbours of $u$ in
$\tilde G$ or in $G$, since the star vertices $v\in N_{\EV}(u)$ are odd in
$\tilde G$, each having lost a single edge, and have $\tilde D\ge 1$ by
Lemma~\ref{int:lem-LB1}(ii)). We obtain $A\subseteq F$ with $|A|\ge\ceil{k/2}$
which can be added back towards $u$, producing a decomposition $D_1$ of
$\tilde G\cup A$ with $|D_1|=|\tilde D|$ and
$D_1(u)\ge \tilde D(u)+\ceil{k/2}\ge 1+\ceil{k/2}$, the bound
$\tilde D(u)\ge1$ coming from Lemma~\ref{int:lem-LB1}(ii) applied to $u$,
which is odd in $\tilde G$.

\emph{Step 4 (reinsertion of the remainder).}
Let $L=F\setminus A$; $|L|=\floor{k/2}=(k-1)/2$. Process the edges of $L$
one at a time. Suppose $j-1$ edges of $L$ have been reinserted, giving a
decomposition $D_j$ of
$\tilde G\cup A\cup(\text{first } j-1 \text{ edges of } L)$ with
\[
D_j(u)\;\ge\;1+\ceil{k/2}-(j-1).
\]
Let $uv\in L$ be the next edge. We apply Lemma~\ref{int:lem-fan6} \emph{towards
$v$}, i.e.\ with the roles of $u$ and $v$ in the statement exchanged: the
condition required is
\[
D_j(u)\;>\;\bigl|\{w\in N_{G_j}(v) : D_j(w)=0\}\bigr|,
\]
where $G_j$ is the current graph, and the operation then yields a
decomposition $D_{j+1}$ of $G_j+uv$ with $D_{j+1}(v)=D_j(v)+1$,
$D_{j+1}(u)=D_j(u)-1$, and $|D_{j+1}|=|D_j|$.

We bound the right-hand side. If $D_j(w)=0$ then $w$ has even degree in
$G_j$ by Lemma~\ref{int:lem-LB1}(iii). We classify the neighbours of $v$ and
their endpoint counts. Odd-degree vertices of $G$ are never incident with
a star edge, keep their degree, and remain odd throughout; they always
have $D\ge1$. A star vertex $v'\in N_{\EV}(u)$, $v'\ne v$, is \emph{not a
neighbour of $v$ at all}: the vertices $v$ and $v'$ are both even-degree
vertices of $G$, so the edge $vv'$ would close the triangle $u,v,v'$ in
$\EV(G)$, contradicting (A1). (This is just as well: a star vertex whose
edge has been reabsorbed towards $u$ by Lemma~\ref{int:lem-fan7} may well have
$D=0$, each absorption step lowering $D(v')$ by one.) The vertex $u$
\emph{is} a neighbour of $v$ and has regained at least $\ceil{k/2}\ge3$
edges, so $D_j(u)\ge3>0$. The only remaining neighbours of $v$ that were
even in $G$ and lie outside $\{u\}\cup N_{\EV}(u)$ are the members of
$N_{\EV}(v)\setminus\{u\}$, of which there are $\dEV(v)-1\le2$. Hence
\[
\bigl|\{w\in N_{G_j}(v) : D_j(w)=0\}\bigr|
\;\le\; \bigl|\{w\in N_{\EV}(v) : w\ne u\}\bigr|
\;=\;\dEV(v)-1\;\le\;2,
\]
using that $v$ is low, i.e.\ $\dEV(v)\le 3$ by (A2). At the last step
$j=(k-1)/2$ we have
$D_j(u)\ge 1+\ceil{k/2}-(j-1)=1+(k+1)/2-(k-3)/2=3>2$, and for earlier steps
the budget is only larger. So each edge of $L$ is addible towards its star
vertex $v$, free of charge.

After all $k$ edges are reinserted we obtain a decomposition of $G$ with
$|D|=|\tilde D|\le\floor{n/2}$, completing the proof.
\end{proof}

The alignment condition is not the only route through a high vertex; the
following theorem needs no alignment hypothesis at all, at the price of a
separation condition between high vertices.

\begin{theorem}[separated high vertices]\label{thm:lasep}
Assume {\rm(A1)}, {\rm(A2)}, and moreover that any two distinct high
vertices of $G$ have no common EV-neighbour. Then $\pn(G)\le\floor{n/2}$.
\end{theorem}
\begin{proof}
We prove the stronger statement, for all graphs in the class connected or
not, by strong induction on $|E(G)|$. If there is no high vertex, every
E-degree is at most $3$ and Theorem~\ref{int:thm-fan} applies component-wise,
so Lemma~\ref{int:lem-sum} gives the bound. Otherwise pick a high vertex $u$
and let $k=\dEV(u)\ge 4$, which is odd by {\rm(A2)}; let $F$ be the E-star
at $u$ and $\tilde G=G-F$. Every star vertex has odd degree in $\tilde G$
(it loses the single edge $uv$), and $u$ has odd degree in $\tilde G$; the
even-degree vertices of $\tilde G$ are precisely the even-degree vertices
of $G$ outside $\{u\}\cup N_{\EV}(u)$, and for such a vertex $x$ its
E-degree in $\tilde G$ equals
$\dEV(x)-|N_{\EV}(x)\cap N_{\EV}(u)|$; when $x$ is high the subtraction
term is $0$, because distinct high vertices have no common EV-neighbour.
Hence $\tilde G$ again satisfies {\rm(A1)}, {\rm(A2)} and the separation
condition, so the induction hypothesis gives a decomposition $\tilde D$ of
$\tilde G$ with $|\tilde D|\le\floor{n/2}$. The neighbours of $u$ in
$\tilde G$ are its odd-degree neighbours in $G$, so each has
$\tilde D\ge 1$ (Lemma~\ref{int:lem-LB1}(ii)); the star vertices are odd in
$\tilde G$, each having lost a single edge, so they too have
$\tilde D\ge1$ by Lemma~\ref{int:lem-LB1}(ii), and the hypothesis of
Lemma~\ref{int:lem-fan7} holds read over $G$; Lemma~\ref{int:lem-fan7} returns an
addible set $A\subseteq F$ with $|A|\ge\ceil{k/2}$ towards $u$, giving
$D_1(u)\ge 1+\ceil{k/2}$. Each edge of $L=F\setminus A$, $|L|=\floor{k/2}$,
is then reinserted towards its star vertex $v$ via Lemma~\ref{int:lem-fan6}:
the vertices $w$ with $D_j(w)=0$ are even in $G_j$, and as in the proof of
Theorem~\ref{thm:claAstar} they are confined to
$N_{\EV}(v)\setminus\{u\}$, of size at most $\dEV(v)-1\le 2$; at the last
step $D_j(u)\ge 3>2$. This rebuilds a decomposition of $G$ with
$\floor{n/2}$ paths.
\end{proof}

\begin{remark}\label{rem:claCO}
The proof of Theorem~\ref{thm:lasep} goes through, verbatim up to one
parity check, with separation replaced by the following weaker condition:
for any two distinct high vertices $x,z$ and any set $S$ of high vertices
disjoint from $\{x,z\}$, the number of common EV-neighbours $w$ of $x$ and
$z$ with $N_{\EV}(w)\cap H\cap S=\emptyset$ is even (equivalently: for
each pair $x,z$, the common EV-neighbours are split by their ``profile''
$\pi(w)=N_{\EV}(w)\cap H$ into classes of even size). Indeed this
condition is again inherited after a star deletion, because a vertex of
$\EV(\tilde G)$ is precisely an even vertex seeing no member of
$\{u\}\cup N_{\EV}(u)$, so the counts for $\tilde G$ with $S'$ are the
original counts for $S=\{u\}\cup S'$.
\end{remark}

Separation and its profile relaxation still leave open the case of high
vertices that do share EV-neighbours. The next theorem settles the case of
exactly two high vertices except for a corner recorded afterwards.

\begin{theorem}[two high vertices]\label{thm:clatwohigh}
Assume {\rm(A1)} and {\rm(A2)}, let $u,x$ be the \emph{only} high vertices
of $G$, and suppose $t_{ux}:=|N_{\EV}(u)\cap N_{\EV}(x)|\ge1$. If either
\begin{enumerate}[label={\rm(\roman*)},leftmargin=2em]
\item some common EV-neighbour $v_1$ of $u$ and $x$ has $\dEV(v_1)=2$, or
\item $x$ has a private EV-neighbour $y^*$, i.e.\
$y^*\in N_{\EV}(x)\setminus N_{\EV}(u)$, with $\dEV(y^*)=1$,
\end{enumerate}
then $\pn(G)\le\floor{n/2}$.
\end{theorem}

\begin{proof}
Let $F$ be the E-star at $u$, $k=\dEV(u)$, and delete one further edge
$e^*$ at $x$: $e^*=xv_1$ in~(i), $e^*=xy^*$ in~(ii), giving
$\tilde G=G-F-e^*$. In $\tilde G$ the vertices $u$ and $x$ are odd, and so
is every former star vertex (each lost the single edge $uv$); in~(ii) the
far endpoint $y^*$ of $e^*$ is odd as well (it lost $xy^*$). The remaining
even vertices --- in~(i) the vertex $v_1$ (which lost both $uv_1$ and
$xv_1$ and stays even) together with the far even vertices and the private
EV-neighbours of $x$; in~(ii) the far even vertices other than $y^*$ and
the private EV-neighbours of $x$ other than $y^*$ --- all have E-degree at
most $3$ in $\tilde G$: the private ones by {\rm(A2)} at $x$, the far ones
because $u$ and $x$ are the only high vertices, and $v_1$ because
$\dEV(v_1)=2$. Hence $\EV(\tilde G)$ is triangle-free, being an induced
subgraph of $\EV(G)$, of maximum E-degree at most $3$, so $\tilde G$ lies
in the class of Theorem~\ref{int:thm-fan} and has, component-wise with
Lemma~\ref{int:lem-sum}, a decomposition $\tilde D$ with
$|\tilde D|\le\floor{n/2}$. In~(i) reinsert $e^*$ towards $v_1$
(Lemma~\ref{int:lem-fan6}: $\tilde D(x)\ge1$ since $x$ is odd, while every
neighbour of $v_1$ in $\tilde G$ is odd except $x$ --- the even-degree
neighbours of $v_1$ in $G$ are exactly $u$ and $x$ --- so the passing
count is $0$); this activates $v_1$. In~(ii) no activation is needed:
every star vertex has lost exactly one edge and is odd. Then
Lemma~\ref{int:lem-fan7} (hypothesis read over $G$, as in the proof of
Theorem~\ref{thm:claAstar}) returns $A\subseteq F$ with $|A|\ge\ceil{k/2}$
addible towards $u$, and each edge of $L=F\setminus A$ is reinserted
towards its star vertex as in the proof of Theorem~\ref{thm:claAstar}
(passing count $\le\dEV(v)-1\le2$, budget $\ge3$ at the last step). In~(ii)
the edge $xy^*$ is reinserted towards $y^*$, whose neighbours except $x$
are all odd, so this step is legal at any position of the sequence. All
$k+1$ deleted edges are restored with $|D|=|\tilde D|\le\floor{n/2}$.
\end{proof}

\begin{corollary}\label{cor:clatwohigh}
If $G$ satisfies {\rm(A1)}--{\rm(A2)} and has at most two high vertices,
then $\pn(G)\le\floor{n/2}$, unless $G$ has exactly two high vertices $u,x$
with $t_{ux}\ge1$, every common EV-neighbour of $u$ and $x$ of E-degree
$3$, no private EV-neighbour of E-degree $1$ on either side, and such
that, for \emph{every} common EV-neighbour $v$ of $u$ and $x$, writing $z$
for the third EV-neighbour of $v$, $F$ for the E-star at $u$ and
$\tilde G=G-F-\{xv\}$, all of the following hold:
$\pn(\tilde G)=\floor{n/2}$; $\beta(\tilde G,x)=\floor{n/2}+1$ (cf.\
Section~\ref{int:subsec-budget}); $z$ is buried (i.e.\ $D(z)=0$) in every
optimal decomposition of $\tilde G$; in every optimal decomposition
$\tilde D$ of $\tilde G$, either $\tilde D(v)=0$, or $\tilde D(v)\ge2$ and
at least $\tilde D(v)$ private EV-neighbours of $x$ have $\tilde D=0$; and
$\tilde G$ has at most one component of odd order.
\end{corollary}
\begin{proof}
$t_{ux}=0$ is Theorem~\ref{thm:lasep} (with only two high vertices the
separation condition reduces to $t_{ux}=0$). The two alternatives are
Theorem~\ref{thm:clatwohigh}. For the remaining graphs fix a common
EV-neighbour $v$, write $z$ for its third EV-neighbour and
$\tilde G=G-F-\{xv\}$; then $\EV(\tilde G)$ again has maximum E-degree
$\le3$, and each of the following five routes closes the case for this
$v$, the star at $u$ being absorbed afterwards exactly as in the proof of
Theorem~\ref{thm:clatwohigh}:
\begin{enumerate}[label={\rm(\alph*)},leftmargin=2em]
\item $\pn(\tilde G)\le\floor{n/2}-1$; take an optimal decomposition
$\tilde D$. If $\tilde D(x)\ge3$ we are on route~(b). Otherwise
$\tilde D(x)=1$ (it is odd by Lemma~\ref{int:lem-LB1}(iii), $x$ being odd in
$\tilde G$), and since $\deg_{\tilde G}(x)=\deg_G(x)-1\ge5$, the unique
path ending at $x$ continues through $x$; cutting it at $x$ adds one path
(at most $\floor{n/2}$ in all) and raises the endpoint count at $x$ to $3$
--- route~(b);
\item a decomposition $\tilde D$ with $|\tilde D|\le\floor{n/2}$ and
$\tilde D(x)\ge3$: then $xv$ is addible towards $v$, the passing count at
$v$ being at most $1$ ($z$ is the only neighbour of $v$ in $\tilde G$ that
can be dead; all other neighbours are odd);
\item such a decomposition with $\tilde D(x)=1$ and $\tilde D(z)\ge1$: the
passing count at $v$ drops to $0$;
\item such a decomposition with $\tilde D(v)\ge2$ larger than the number of
dead private EV-neighbours of $x$: reinsert $xv$ towards $x$ (the passing
count at $x$ sees only private EV-neighbours of $x$, the other neighbours
of $x$ being active), so that $v$ remains active and the full star,
including $uv$, is absorbable;
\item $\tilde G$ has at least two components of odd order: the
component-wise Theorem~\ref{int:thm-fan}, summed with Lemma~\ref{int:lem-sum},
gives route~(a).
\end{enumerate}
The budget arithmetic of the reinsertion skeleton behind routes (a)--(d),
and the formal statement that the exception of this corollary is exactly
the simultaneous failure of all five routes at \emph{every} common
EV-neighbour, are given in Appendix~\ref{com:app-band}.
The stated exception is exactly the case in which, for every common
EV-neighbour $v$, none of the five routes applies.

\emph{Band closure of the corner.} One further narrowing is available in
the constructed family of Remark~\ref{rem:clagap}. Route~(a) can fail
there only when $\tilde G$ carries a \emph{deadlock band} (an even vertex
of $\tilde G$ with more corridor neighbours than terminal units; see
Remark~\ref{rem:clagap} for the definition), because a band-free
$\tilde G$ in the family has
$\pn(\tilde G)=\mathrm{odd}(\tilde G)/2$; and the counting theorem forces
$\tilde G$ to be band-free whenever the counting gap
$\floor{n/2}-\mathrm{odd}(\tilde G)/2$ --- which is at least $1$ in this
situation --- is at most $2$, on the family proper (the constructed family
without cross-private pairs). Hence for every family-proper instance with
counting gap at most $2$, route~(a) applies unconditionally: the corner is
empty at theorem level, not merely computationally; the cross-private
instances are backed at probe level by the machine certificates of
Appendix~\ref{com:app-comp}. For larger gaps the corner is blocked only by the
excess bound recorded in Remark~\ref{rem:clagap}.
\end{proof}

\begin{corollary}\label{cor:clauniquehigh}
If $G$ satisfies (A1) and (A2) and has \emph{exactly one} high vertex,
then $\pn(G)\le\floor{n/2}$.
\end{corollary}

\begin{proof}
Let $u$ be the unique high vertex. Every vertex other than $u$ is low, so
every even vertex $x\notin\{u\}\cup N_{\EV}(u)$ satisfies
$\dEV(x)-|N_{\EV}(x)\cap N_{\EV}(u)|\le\dEV(x)\le 3$; this is condition
(A4) of Theorem~\ref{thm:claAstar}. (Alternatively, with only one high
vertex the separation hypothesis of Theorem~\ref{thm:lasep} holds
vacuously.)
\end{proof}

We now formulate the statement with neither the alignment nor the
separation condition, which is the form investigated in the computational
study behind this paper and which all our computations support.

\begin{conjecture}\label{conj:claAstar}
Let $G$ be a connected graph such that $\EV(G)$ is triangle-free, and
every vertex $v$ with $\dEV(v)\ge 4$ has odd $\dEV(v)$ and only
EV-neighbours $w$ with $\dEV(w)\le 3$. Then $\pn(G)\le\floor{n/2}$.
\end{conjecture}

\begin{remark}[where the general argument stops]\label{rem:clagap}
The class {\rm(A1)--(A2)} is not preserved by the only deletion the engine
supports. Deleting the E-star at a high vertex $u$ turns every star vertex
odd; an even vertex $x$ outside the star keeps its degree but loses the
E-degrees contributed by star vertices, and its E-degree in $\tilde G=G-F$
equals $\dEV(x)-t_x$ with $t_x=|N_{\EV}(x)\cap N_{\EV}(u)|$. If $x$ is a
high vertex with $t_x$ odd and $\dEV(x)-t_x\ge 4$, then $x$ sits in
$\EV(\tilde G)$ at an even E-degree, the class is lost, and neither
Theorem~\ref{int:thm-fan} nor the induction applies to $\tilde G$. Concretely,
in the double star on
$\{u,x,v_1,w_1,\dots,w_4,y_1,\dots,y_4\}$ with edges
$uv_1,uw_1,\dots,uw_4,xv_1,xy_1,\dots,xy_4$ and pendant leaves attaching
so that the eleven listed vertices have even degree, both $u$ and $x$ are
high with $\dEV=5$, $t_{ux}=1$ is odd, and after deleting the star at $u$
the vertex $x$ has E-degree $4$, even. This graph (and its even-order
one-leaf extension, on which our exact programs give $\pn=5\le 11$)
satisfies {\rm(A1)--(A2)} but none of {\rm(A4)}, the separation condition,
or the profile condition of Remark~\ref{rem:claCO}; a minimal
counterexample to Conjecture~\ref{conj:claAstar}, if it exists, must have
a high vertex with such a bad partner at {\rm every} high vertex
(Proposition~\ref{prop:claresidue}). No counterexample is known: exact
dynamic programs verify $\pn\le\floor{n/2}$ on all tested instances
satisfying {\rm(A1)--(A2)} ($166$ instances in the computational study
behind this paper, a further $413$ in Appendix~\ref{com:app-comp}, zero
violations), and the family above sits far from the bound
($\pn=\mathrm{odd}(G)/2$).

The ``one endpoint unit'' obstruction recorded in an earlier draft of this
material was an artefact: the reinsertion bookkeeping it refers to
restores $k+2$ edges after deleting $k+1$, i.e.\ counts the patched edge
$xv_1$ twice (an activation step towards $v_1$ and a final step towards
$x$). Once $x$'s endpoint loss is recognised as cost-free --- $|D|$ is the
only invariant, and the argument of Theorem~\ref{thm:clatwohigh} involves
no induction --- the obstruction disappears for all two-high graphs except
the corner recorded in Corollary~\ref{cor:clatwohigh}.

The exceptional graphs of Corollary~\ref{cor:clatwohigh} are subject to
five simultaneous constraints, of which only the first two were recorded
when the corner was first isolated: the burial of the third EV-neighbour
$z$ and of $v$ (or of enough private neighbours of $x$) in \emph{every}
optimal decomposition of $\tilde G$ (with $\tilde G$ as in
Corollary~\ref{cor:clatwohigh}), and the parity of the components of
$\tilde G$, are each individually sufficient to reopen the reinsertion
argument (the edge $xv$ can then be added towards $v$ at zero passing
cost, or a second deletion $G-\tilde G$ becomes unnecessary). No instance
satisfying even the original two constraints was found: across $207$
constructed instances of the dirty class---including the entire range
where the counting gap $\floor{n/2}-\mathrm{odd}(\tilde G)/2$ attains its
minimum $1$, which the degree capacity of the two far vertices forces to a
single parameter-free shape---a decomposition of $\tilde G$ with at most
$\floor{n/2}-1$ paths was always produced, so the corner never materializes
in computation.

The band-free tightness statement is now a theorem for the \emph{constructed
family} (the two-high graphs satisfying the invariants of our construction
pipeline: no $\EV$-edge joins two far even vertices of $\tilde G$; the
odd-neighbour blocks of $u$ and of $x$ are of leaf type (an odd number of
pendant leaves) or gadget type, possibly mixed; every far vertex of
$\tilde G$ has degree at most $4$; the $207$ instances above and the $125$
decision-level instances of Appendix~\ref{com:app-comp} all lie in it).
Writing, for an even vertex $f$ of $\tilde G$, $Y(f)$ for its
\emph{corridor} neighbours --- the degree-two even vertices adjacent to
both $x$ and $f$ --- and $N(f)$ for its neighbour set in $\tilde G$, call
$f$ a \emph{deadlock band} if it has more corridor neighbours than
terminal units, where the terminal units of $f$ are its remaining incident
edges (in symbols, $|Y(f)|>|N(f)|-|Y(f)|$). If no even vertex is a
deadlock band, then $\tilde G$ admits a \emph{paired} decomposition:
exactly $\mathrm{odd}(\tilde G)/2$ paths with all and only the odd
vertices of $\tilde G$ as endpoints. The proof is a five-step pairing
scheme---self-contained odd-neighbour blocks, forced fusions of the
corridors and of the unit at the common neighbour, an injective pairing of
corridor ends with distinct terminal units at every far vertex (which is
precisely the negation of the band condition), and a sorted offset
matching at $x$ that never joins two corridors of the same far vertex,
possible because a far vertex of the family has degree at most $4$ and
hence at most two corridors. Every resulting trail visits $x$ once and
leaves each far vertex for good, so no trail closes. A single band, on
the other hand, costs exactly one path: spending the free edge of $x$ on a
surplus corridor makes that corridor a double endpoint and lets its
partner ride through $x$ to a terminal, matching the exact value
$\pn(\tilde G)=\mathrm{odd}(\tilde G)/2+1$ computed by dynamic program.
Consequently the corner of Corollary~\ref{cor:clatwohigh} is empty
whenever the counting gap is at most $2$, unconditionally on the
constructed family proper, with the cross-private extension covered at
probe level by the machine certificates of Appendix~\ref{com:app-comp}; for
larger gaps it remains blocked only by the excess bound
$\mathrm{excess}\le\lfloor y/2\rfloor$, where
$\mathrm{excess}:=\pn(\tilde G)-\mathrm{odd}(\tilde G)/2$ and $y$ is the
number of corridors --- verified on all $235$ computed instances of the
family, the provable bound being
$\pn(\tilde G)\le \mathrm{odd}(\tilde G)/2+2B-1$ for $B$ bands, with a
single band forcing exactly $+1$. The restriction to the constructed
family is now known to be sharper than necessary in two directions and
irreducible in a third. Far--far edges of the even subgraph and far
vertices carrying several parity leaves, both excluded by the
construction, can be admitted: the corridor ends at a far vertex $f$
admit an injective pairing with terminal units and far--far ends whenever
$|Y(f)|\le t(f)+d(f)$ (where $t(f)$ is the number of terminal units of
$f$ and $d(f)$ its far--far degree), and the matching at $x$ survives
arbitrary corridor multiplicities because a far vertex of the class has at
most three even neighbours, so a rearrangement argument replaces the
sorted offset. What cannot be admitted are arbitrary odd-neighbour
blocks: a block whose component has exactly two odd vertices admits a
paired decomposition if and only if it admits a Hamilton path between
them, and the cycle-with-one-pendant and double-triangle blocks have
none, so $\tilde G$ can be untight while carrying no band at all---the
band-free tightness statement is a phenomenon of the constructed family,
not of the class, and these cycle blocks are the first known
counterexamples to the general band-free statement. (Same-side pairs of
private neighbours, a third candidate obstruction, are excluded by the
triangle-free condition {\rm(A1)} already.) The counting gap of the dirty
class also assumes that every even vertex lies in the skeleton; blocks
with even interior vertices shift $\floor{n/2}$ and must be recomputed
directly (see Appendix~\ref{com:app-comp} for the machine certificates).
\end{remark}

\begin{proposition}\label{prop:claresidue}
Let $G$ be a minimal counterexample to Conjecture~\ref{conj:claAstar}.
Then every high vertex $u$ of $G$ admits a far high vertex $x$ with
$|N_{\EV}(x)\cap N_{\EV}(u)|$ odd and
$\dEV(x)-|N_{\EV}(x)\cap N_{\EV}(u)|\ge4$. In particular, by
Theorem~\ref{thm:clatwohigh} a minimal counterexample has at least three
high vertices, unless it is one of the two-high exceptional graphs of
Corollary~\ref{cor:clatwohigh} (empty on the constructed family proper
whenever the counting gap is at most two, by the band-free tightness
material of Remark~\ref{rem:clagap}; no instance was found in our
computations, in or beyond the family).
\end{proposition}
\begin{proof}
Read ``minimal'' as minimal with respect to the number of edges, and let
$u$ be a high vertex of $G$. We show that if $u$ has no bad partner, then
$\pn(G)\le\floor{n/2}$, contradicting the choice of $G$; hence $u$ has a
bad partner. Suppose, then, that no far even vertex $x$ satisfies $t_x$
odd and $\dEV(x)-t_x\ge4$, where $t_x=|N_{\EV}(x)\cap N_{\EV}(u)|$. Let
$F$ be the E-star at $u$ and $\tilde G=G-F$. As in Step~1 of the proof of
Theorem~\ref{thm:claAstar}, the even-degree vertices of $\tilde G$ are
precisely the far even vertices of $G$, $\EV(\tilde G)$ is an induced
subgraph of $\EV(G)$ and hence triangle-free, and a far vertex $x$ has
E-degree $\dEV(x)-t_x$ in $\tilde G$. Thus every high vertex of
$\tilde G$ is a far even vertex $x$ of $G$ with $\dEV(x)-t_x\ge4$; by our
assumption $t_x$ is even for such an $x$, and since $\dEV(x)\ge4$ makes
$x$ high in $G$, (A2) gives $\dEV(x)$ odd, so the E-degree
$\dEV(x)-t_x$ of $x$ in $\tilde G$ is odd. Moreover the EV-neighbours $w$
in $\tilde G$ of such an $x$ satisfy
$\dEV(w)-t_w\le\dEV(w)\le3$, the last inequality by (A2) applied to the
high vertex $x$ in $G$. Hence $\tilde G$ satisfies {\rm(A1)} and
{\rm(A2)}. Each component of $\tilde G$ is a connected graph satisfying
{\rm(A1)}--{\rm(A2)} with fewer edges than $G$ (as $u$ is high,
$F\neq\emptyset$), so by minimality each admits a decomposition with at
most $\floor{n_c/2}$ paths, and Lemma~\ref{int:lem-sum} gives a decomposition
$\tilde D$ of $\tilde G$ with $|\tilde D|\le\floor{n/2}$. Now the
reinsertion argument of Steps~3 and~4 of the proof of
Theorem~\ref{thm:claAstar} uses (A4) only through the existence of such a
$\tilde D$: the star vertices are odd in $\tilde G$, so the hypothesis of
Lemma~\ref{int:lem-fan7} holds read over $G$ exactly as in Step~3, and each
edge of $F\setminus A$ is addible towards its star vertex against a
passing count of at most $\dEV(v)-1\le2$ (by (A2)) and a budget of at
least $3$. The full star is restored with $|D|=|\tilde D|\le\floor{n/2}$
paths --- a contradiction.
\end{proof}

We also record a structural consequence of (A1) on odd orders, used below
to place the sharpness family of Proposition~\ref{prop:claoscsharp}
relative to the class.

\begin{proposition}\label{prop:claodddensity}
Let $G$ be connected of odd order $n$ with $\EV(G)$ triangle-free. Then
$m\le (n-1)^2/2$; equivalently, no odd semi-clique has triangle-free
$\EV(G)$.
\end{proposition}
\begin{proof}
Suppose $c=\binom{n}{2}-m\le (n-3)/2$. Let $S$ be the set of even-degree
vertices, $s=|S|$; since $n$ is odd, a vertex has odd degree iff its
complement-degree is odd, so each vertex outside $S$ has positive
complement-degree. Let $t$ be the number of complement edges inside $S$.
Since $G[S]=\EV(G)$ is triangle-free, Mantel's theorem gives
$\binom{s}{2}-t\le\floor{s^2/4}$, i.e.\ $t\ge \binom{s}{2}-\floor{s^2/4}$.
Counting complement degrees over $S$ and over its complement separately,
$2c\ge 2t+(n-s)$, so $2t\le s-3$. If $s$ is even this reads
$s(s-2)/2\le s-3$, i.e.\ $s^2-4s+6\le0$; if $s$ is odd,
$(s-1)^2/2\le s-3$, i.e.\ $s^2-4s+7\le0$; both are impossible. (Note
$s\ge1$: the number of odd-degree vertices is even while $n$ is odd.)
\end{proof}

We record the sharpness of the triangle-free requirement. By
Proposition~\ref{prop:claodddensity} this failure is a density phenomenon:
on odd orders, high density and (A1) are mutually exclusive, so the
following family is forced.

\begin{proposition}\label{prop:claoscsharp}
For every $k\ge 2$, the graph $G_k=K_{2k+1}-(k-1)K_2$ (delete a matching
of size $k-1$) satisfies $\EV(G_k)\cong K_3$, hence violates (A1), and
$\pn(G_k)\ge k+1=\ceil{n/2}>\floor{n/2}$. Moreover $\pn(G_k)=k+1$ for the
tested values $k\in\{2,3,4,5,6,8,10\}$, with explicit decompositions
machine-verified.
\end{proposition}

\begin{proof}
The $k-1$ matched vertices have degree $2k-1$ and the three unmatched
vertices have degree $2k$; the latter are pairwise adjacent, so
$\EV(G_k)$ is a triangle on them. The edge count is
$m=\binom{2k+1}{2}-(k-1)=2k^2+1$, so Lemma~\ref{int:lem-LB2} gives
$\pn\ge\ceil{(2k^2+1)/2k}=k+1$; equivalently $G_k$ is an odd semi-clique
and Observation~\ref{int:obs-density} applies. This proves the necessity
claim for all $k$. For the tested values, explicit decompositions into
$k+1$ paths (obtained from the Walecki cycle decomposition of $K_{2k+1}$,
spreading the matching edges over distinct cycles, cutting each affected
cycle at its matching edge and splitting the untouched cycle into a path
plus a single edge) were verified by machine to partition $E(G_k)$
exactly; for $k=2$ the value $\pn(K_5-M_1)=3$ is also certified by exact
dynamic programming.
\end{proof}

\begin{remark}
The exactness claim in Proposition~\ref{prop:claoscsharp} is stated only
for tested $k$: the existence, for all $k$, of a Walecki-type cycle
decomposition of $K_{2k+1}$ in which a given $(k-1)$-matching is spread
over distinct cycles is plausible but we do not prove it.
\end{remark}

Next, the witness family $W(m)$, which motivates Theorem~\ref{thm:claAstar}:
it lies in no published class we are aware of, yet satisfies all
hypotheses.

\begin{figure}[h]
\centering
\begin{minipage}{0.92\textwidth}\small
For odd $m\ge 5$, let $W(m)$ have vertex set
$\{a_1,a_2\}\cup\{b_1,\dots,b_m\}\cup\{\ell_1,\ell_2\}\cup\{h,h'\}\cup\{x_1,\dots,x_6\}$
and edge set
\[
E(K_{2,m}) \;\cup\; \{a_1\ell_1,\,a_2\ell_2\} \;\cup\;
\{b_1h,\,b_1h'\} \;\cup\; \{hx_1,\,hx_2,\,h'x_1,\,h'x_2\} \;\cup\; E(K_6),
\]
where $K_6$ has vertex set $\{x_1,\dots,x_6\}$. Precisely: $h$ and $h'$
each have neighbourhood $\{b_1,x_1,x_2\}$.
\end{minipage}
\caption{The witness family $W(m)$.}\label{fig:claWm}
\end{figure}

\begin{lemma}\label{lem:claWmEV}
For every odd $m\ge 5$, the graph $W(m)$ of Figure~\ref{fig:claWm}
satisfies $\EV(W(m))\cong K_{2,m}$ exactly, hence satisfies (A1) and (A2)
with $\dEV(a_1)=\dEV(a_2)=m$ and $\dEV(b_j)=2$ for all $j$.
\end{lemma}

\begin{proof}
The degrees are:
$\deg(a_i)=m+1$, even; $\deg(b_j)=2$ for $j\ge2$ and $\deg(b_1)=4$, all
even; $\deg(\ell_i)=1$, $\deg(h)=\deg(h')=3$, and $\deg(x_i)=5+2=7$ for
$i\in\{1,2\}$, $\deg(x_i)=5$ for $i\in\{3,\dots,6\}$, all odd. Hence the
even-degree vertex set is exactly $\{a_1,a_2\}\cup\{b_1,\dots,b_m\}$, and
the induced subgraph is $K_{2,m}$: triangle-free, with $\dEV(a_i)=m$ odd
and all EV-neighbours of $a_i$ of E-degree $2\le3$. Note that the two
degree-$3$ vertices $h,h'$ attach to $b_1$ \emph{in pairs} (both $h$ and
$h'$), keeping $b_1$ even; attaching a single such vertex would flip $b_1$
to odd degree and destroy the structure.
\end{proof}

\begin{proposition}\label{prop:claWmbound}
For every odd $m\ge 5$, $\pn(W(m))\le\floor{n/2}$ with $n=m+12$.
\end{proposition}

\begin{proof}
By Lemma~\ref{lem:claWmEV} the graph satisfies (A1), (A2). We verify (A4)
at $u=a_1$: the only even-degree vertex outside
$\{a_1\}\cup\{b_1,\dots,b_m\}$ is $a_2$, and, since
$N_{\EV}(a_1)=N_{\EV}(a_2)=\{b_1,\dots,b_m\}$,
$\dEV(a_2)-|N_{\EV}(a_2)\cap N_{\EV}(a_1)|=m-m=0\le3$; every vertex
outside this set other than $a_2$ has odd degree, hence is not in
$\EV(\tilde G)$ after the deletion and is irrelevant (its E-degree in
$\tilde G$ is $0$). Thus (A4) holds, and Theorem~\ref{thm:claAstar} gives
$\pn(W(m))\le\floor{(m+12)/2}$. For $m\in\{5,7,9\}$ explicit
decompositions into exactly $\floor{n/2}=8,9,10$ paths were in addition
machine-verified.
\end{proof}

We now compare Theorem~\ref{thm:claAstar} with the published record. The
following two statements make the incomparability precise; both witnesses
have exactly computed path numbers.

\begin{proposition}[escape]\label{prop:claescape}
Let $m\ge5$ be odd and let $W(m)$ be as in Figure~\ref{fig:claWm}. Then
$\EV(W(m))\cong K_{2,m}$, whence Theorem~\ref{thm:claAstar} gives
$\pn(W(m))\le\floor{n/2}$ with $n=m+12$. No published theorem known to us
implies this bound:
\begin{itemize}[leftmargin=1.5em]
\item every member of the Fan line \cite{Fan05,BS21,FHZ23} requires every
block of $\EV(G)$ to have maximum degree at most $3$; the graph
$\EV(W(m))\cong K_{2,m}$ is $2$-connected, hence a single block of maximum
degree $m\ge5>3$, so this requirement fails (the block level, not the
global value of $\Delta(\EV)$, is the discriminating parameter ---
compare $K_{1,5}$ with one pendant leaf at each of its six vertices, where
$\Delta(\EV)=5$ while every block of $\EV$ is a single edge);
\item $W(m)$ contains a $K_6$, so $\mathrm{tw}(W(m))\ge5$ and its
degeneracy is at least $5$, which excludes \cite{LY24} (treewidth $\le4$),
\cite{BCLS20} (treewidth $\le3$) and \cite{ZLH24} ($3$-degenerate);
\item $\Delta(W(m))=\max(m+1,7)\ge7$, which excludes \cite{BP19}
($\Delta\le5$) and \cite{CFL} ($\Delta\le6$ with an independence
condition);
\item $\Delta(W(m))\ge7$, so $W(m)$ is not subcubic and the
polynomial-time algorithm of \cite{MY26} for subcubic graphs does not
apply;
\item $\EV(W(m))\cong K_{2,m}$ contains a $4$-cycle, so it is not a forest
and \cite{Pyb96} does not apply;
\item $W(m)$ is not Eulerian, so \cite{FK88} does not apply; it is not a
block graph, so \cite{CW25} does not apply.
\end{itemize}
For $m\in\{5,7,9\}$ every membership and non-membership claim above,
together with explicit decompositions into $\floor{n/2}$ paths, is
machine-certified.
\end{proposition}

\begin{proposition}[proper inclusion in both directions]\label{prop:claincomp}
The class of Theorem~\ref{thm:claAstar} neither contains nor is contained
in the Botler--Sambinelli class $\mathcal{G}$ \cite{BS21} or the
Fan--Hou--Zhou class \cite{FHZ23}.
\begin{enumerate}[label={\rm(\roman*)},leftmargin=2em]
\item Let $G$ be $K_{2,5}$ with one pendant leaf attached to each of the
two degree-$5$ vertices. Then $\EV(G)\cong K_{2,5}$ is triangle-free, the
two degree-$5$ vertices have odd E-degree $5$ with all E-neighbours of
E-degree $2$, so $G$ satisfies (A1)--(A2); but the graph
$\EV(G)\cong K_{2,5}$ is $2$-connected, hence a single block of maximum
degree $5>3$, so $G$ lies in none of \cite{Fan05,BS21,FHZ23}. Moreover
$\pn(G)=3\le\floor{9/2}$ by exact computation. (For the record, $G$ is
$2$-degenerate, so its bound is also implied by \cite{ZLH24}; this witness
is used only for the incomparability with the Fan line.)
\item Let $G'$ be a triangle $\{t_1,t_2,t_3\}$ together with two vertices
$w_1,w_2$ adjacent only to $t_1$ and to three pendant leaves each, so that
$\EV(G')$ consists of the triangle $\{t_1,t_2,t_3\}$ plus the two bridge
edges $t_1w_1$, $t_1w_2$. Every block of $\EV(G')$ has maximum degree at
most $3$ and its single component contains exactly one block with
triangles, so $G'$ lies in $\mathcal{G}$ and in the Fan--Hou--Zhou class;
yet the triangle in $\EV(G')$ violates (A1), so Theorem~\ref{thm:claAstar}
is silent on $G'$. Note $G'$ is not one of the degenerate exceptions of
the ``Gallai-or-SET'' alternative of \cite{BS21}. Exactly,
$\pn(G')=4\le\floor{11/2}$ by exact computation, while the main theorem of
\cite{BS21} gives only $\pn(G')\le\ceil{11/2}=6$ on its class (the
Botler--Sambinelli bound is stated with Gallai's ceiling, as noted in
Section~\ref{sec:intro}); the Fan--Hou--Zhou membership of $G'$ rests on
the block structure of $\EV(G')$ alone and is unaffected.
\end{enumerate}
\end{proposition}

\begin{remark}[formalization and provenance status of the Fan-type chain]
\label{rem:claformal}
The Lean formalization of the statements of this subsection is in
progress: in the current Lean sources the declarations corresponding to
Theorem~\ref{thm:claAstar}, Theorem~\ref{thm:lasep} and their supporting
reinsertion lemmas (\texttt{astar}, \texttt{astar\_separated},
\texttt{fan6\_use}, \texttt{fan7\_use}, \texttt{fan\_class}) carry
admitted (``\texttt{sorry}'') placeholders --- five in total --- so none
of them may be quoted as machine-verified or formally proved, and the
paper cites the chain as paper proofs with formalization in progress.
One research-technical lemma of the formalization archive
(\texttt{fan6\_nf}, a pointwise-profile strengthening of
Lemma~\ref{int:lem-fan6}) was refuted by a machine counterexample, is sealed
as a research record, and is never cited in this paper; the Lean
formalization of Lemma~\ref{int:lem-fan6} uses the four-clause form obtained
by deleting the final pointwise clause $D(w)=D'(w)$ for the remaining
vertices, and no application in this paper consumes that clause. On the
provenance side: the statements of Lemma~\ref{int:lem-fan6} and
Lemma~\ref{int:lem-fan7} follow the restatement in \cite{BS21} (Lemmas 6 and
7 there); the proofs in \cite{Fan05} could not be consulted and the
transcription has not been independently checked against the original.
The machine evidence is direction-corrected: the bidirectional rescan of
Lemma~\ref{int:lem-fan6} is complete at $n\le6$, all
$1{,}701{,}141{,}150$ ordered instances at $n=6$ (the fixed-orientation
half exactly reproducing the historical census of $850{,}570{,}575$), with
zero violations.
\end{remark}

\subsection{The placeability classification programme on
\texorpdfstring{$K_{2k+1}$}{K(2k+1)}}\label{subsec:echar}

\subsubsection{The layer domain}\label{subsubsec:layerdomain}

Fix the layer domain: the carrier is $K_{2k+1}$ on
$V=\mathbb{Z}_{2k}\cup\{\infty\}$; a \emph{type} is an isomorphism class of
$k$-edge loopless graphs without isolated vertices; and for such a type $F$
we say that $F$ is \emph{sharp at layer $k$} ($\mathrm{sharp}_k(F)$) if
there is an injection $\phi:V(F)\to V$ whose $k$ image-edge classes hit
$\mathbb{Z}_k$ exactly, where the edge classes are given by the closed
form
\begin{equation}\label{eq:claclass}
\mathrm{class}(\{u,v\})=\bigl\lfloor(u+v)/2\bigr\rfloor \bmod k,
\qquad
\mathrm{class}(\{\infty,y\})=y \bmod k .
\end{equation}
A placement realizing $\mathrm{sharp}_k(F)$ is a \emph{witness}. Write
\[
\mathcal{E}_k:=\{F:\ \neg\,\mathrm{sharp}_k(F)\}
\]
for the \emph{exception set} at layer $k$, and $K_n@k$ for the sharpness
question for $K_n$ at layer $k$. The values $k=\binom n2$ are the
\emph{triangular layers}; $K_3@3$, $K_4@6$ and $K_5@10$ are sharp (three
positive controls), so $n\ge6$ is the exact lower bound for exceptions
among complete graphs.

\subsubsection{The first four complete-graph layers}\label{subsubsec:firstlayers}

The exact status of the first four critical layers is as follows; tiers
are quoted verbatim and travel with every use of the numbers.

\begin{enumerate}[leftmargin=2.2em,itemsep=0.2em]
\item \textbf{$K_6@15$ (canonical).} $K_6$ is not sharp at $k=15$, and
$\mathcal{E}_{15}=\{K_6\}$ exactly --- the only layer at which the
uniqueness side is closed. The evidence is four-fold: three independent
depth-first implementations exhaust the same $145{,}923$-node search tree
without witness and without cap, and a direct combinatorial exhaustion of
all $\binom{31}{6}=736{,}281$ six-subsets of the carrier (the image-edge
set of a $K_6$ placement depends only on the six chosen vertices) finds
zero placements with $15$ distinct classes.
\item \textbf{$K_7@21$ (machine universal-closure candidate, not
independently audited).} $K_7$ is not sharp at $k=21$, closed on three
independent lines: an engine exhaustion of $2{,}868{,}963$ nodes without
witness and without cap; a regime census of $3{,}071{,}958$ configurations;
and an infinity-census of $3{,}472{,}896$ configurations (all zero hits).
\item \textbf{$K_8@28$ (machine universal-closure candidate, not
independently audited).} $K_8$ is not sharp at $k=28$: a nine-line census
over $354{,}484{,}845$ configurations with zero hits, zero truncation, and
no sieve dependence (the sieve-excluded regimes were also actually
scanned and are empty).
\item \textbf{$K_9@36$ (machine universal-closure candidate, not
independently audited).} $K_9$ is not sharp at $k=36$: a configuration-level
true enumeration over $18{,}385{,}772{,}680$ configurations (nine regime
lines summing to $\binom{71}{8}=10{,}639{,}125{,}640$ --- the Vandermonde
identity is machine-checked --- plus the infinity line
$\binom{36}{8}\cdot2^8=7{,}746{,}647{,}040$, covered in four shards with an
exact non-overlapping-partition assertion), every line zero hits,
doubly cross-validated against the layered-arithmetic verdict below.
\end{enumerate}

\begin{remark}[a registered prediction that must not be read as a
measurement]\label{rem:claregistered}
An earlier registration predicted for $K_8@28$ ``surviving regimes only
$\{q=0,q=7\}$''; the phrase is a registered prediction, not a measured
residual. The actual census found zero surviving regimes on every line,
and the registered phrase must not be quoted as the census outcome.
\end{remark}

\subsubsection{The characterization candidate and the propagation
bound}\label{subsubsec:echarcand}

The four layers above are exactly what the following candidate predicts.
In its literal form: for every $k\ge1$,
\begin{equation}\label{eq:claechar}
\mathcal{E}_k \;=\; \{\,K_n\ :\ n\ge6,\ \tbinom n2=k\,\},
\end{equation}
equivalently the three-part decomposition: (i) $K_n@\binom n2$ is not
sharp for all $n\ge6$ (sharp at $n=3,4,5$, making $n\ge6$ the exact lower
bound); (ii) $\mathcal{E}_k$ is empty off the triangular layers; (iii) on
triangular layers the exception is exactly the complete graph.

\begin{remark}[registry-layer status of \eqref{eq:claechar}; travels with
any citation]\label{rem:claregistry}
Three registration-layer facts bind every use of \eqref{eq:claechar}.
\emph{First}, two competing candidate shapes are on record --- the literal
shape \eqref{eq:claechar} and the structural shape
``$\mathcal{E}_k$ = the types containing a $K_v$ component'' --- and the
two have never been adjudicated against each other. \emph{Second}, the
literal uniqueness side (iii) stands in tension with the propagation
theorem (a super-type is sharp only if each of its components is
placeable; contrapositively, an unplaceable component makes every
\emph{same-layer} type containing it unsharp): since $K_7$ is unsharp at
layer $21$, at $k=21$ the literal shape would force
$|\mathcal{E}_{21}|=1$, while the structural shape predicts a larger set.
An earlier registration claimed the stronger quantitative form
$|\mathcal{E}_{21}|\ge740{,}226$ via the map $T\mapsto K_7\sqcup T$ on the
$14$-edge types; that claim was withdrawn on review --- $K_7\sqcup T$ has
$21+14=35$ edges and is \emph{not} a layer-$21$ type, so the construction
has no objects in its claimed quantifier domain (and cross-layer
inheritance of unplaceability has known counterexamples in the positive
direction). The honest paper-level reading is
$|\mathcal{E}_{21}|\ge1$ (the type $K_7$ itself, under both shapes), and
uniqueness of the minimal core survives as an open question for $k\ge21$.
\emph{Third}, the candidate
carries mixed candidate/profile load and is not independently audited;
its closed faces are: side (ii) to $k\le16$ (including the $k=17$
full-type sweep of $53{,}394{,}755$ types, every type sharp, tier
machine universal-closure candidate, not independently audited), side (i)
to $n=13$ (tower plus census at $n=6..9$ and the layered-arithmetic
deaths of $K_{11}@55$, $K_{12}@66$, $K_{13}@78$), and side (iii) only at
$k=15$; cumulative over $k=2..17$, $69{,}698{,}732$ types were scanned
with exactly one counterexample ($K_6@15$). Open faces: layers $k\ge18$
(other than $17$ itself), tower layers $n\ge14$, uniqueness of the
minimal core for $k\ge21$, and the prime-family infinity segment
(Section~\ref{subsubsec:primefam}).
\end{remark}

\subsubsection{The prime-family theorem at its non-infinity scope}
\label{subsubsec:primefam}

\begin{theorem}[prime family, non-infinity subdomain]\label{thm:claprimefam}
Let $n=p+1$ with $p\ge7$ a prime congruent to $3 \bmod 4$, and let
$k=\binom n2$ (which is even, divisible by $p$, with $p\mid n-1$). Then
every \emph{non-infinity} regime of $K_n@k$ is dead on paper: the
arithmetic sieve forces $q\equiv eo\equiv0\pmod p$, leaving exactly the
allowed regimes $\{(n,0),(1,p),(p,1)\}$, and each of the three dies by
the order-$2$ criterion, whose feasibility equation reduces to ``$n$ must
be a perfect square'' --- contradicting $p\ge7$.
\end{theorem}

\begin{proof}
Write $e$ and $o$ for the numbers of even- and odd-degree placed vertices,
$q=eo$ for the number of odd--odd pairs, and $\Sigma u$ for the sum of the
half-values $u_i=(y_i\pm k)/2$-type representatives of the finite placed
vertices. Two identities hold for any placement: the total class sum
reads $k(k-1)/2+k\cdot t$ with $t$ the number of half-sums in the upper
half-range, and twice the sum of the half-sums over all pairs equals
$(n-1)\Sigma u-q$ (each odd--odd pair loses one in the halving). Since the
class multiset must be exactly $\mathbb{Z}_k$, summing the classes gives
$(n-1)\Sigma u=k(k-1)+2kt+q$. Reducing mod a prime $p$ with $p\mid k$ and
$p\mid(n-1)$ annihilates the first two terms, forcing $q\equiv0\pmod p$. With
$e+o=n-1=p+1$ we get $q=eo\equiv e(1-e)\pmod p$; since $p$ is prime,
$p\mid e$ or $p\mid e-1$, and $1\le e\le p$ forces $e\in\{1,p\}$ (the
regimes $(1,p),(p,1)$; the case $e=p+1$ i.e.\ $o=0$ is the regime
$(n,0)$, where $q=0$ trivially). These are exactly the three allowed
regimes.

Now $k$ is even, so the order-$2$ (second-character) criterion applies:
writing $x_E=\sum(-1)^{u_i}$ over even points and
$x_O=\sum(-1)^{u_j}$ over odd points, feasibility requires
$x_E^2+2x_Ex_O-x_O^2=e-o$, equivalently
$2x_E^2-(x_E-x_O)^2=e-o$, with $|x_E|\le e$, $|x_O|\le o$,
$x_E\equiv e\pmod2$, $x_O\equiv o\pmod2$.
\emph{Regime $(n,0)$}: $o=0$, so $x_O=0$ and $x_E^2=n$: $n$ must be a
perfect square. \emph{Regime $(1,p)$}: $x_E=1$ (from $x_E\equiv1$,
$|x_E|\le1$), so $(x_O-1)^2=p+1=n$: again $n$ must be a perfect square.
\emph{Regime $(p,1)$}: $x_O=\pm1$, so $(x_E+x_O)^2=n$: the same
conclusion. For $p\ge7$, $n=p+1$ lies strictly between $3^2=9$'s
neighbours $(\lceil\sqrt{8}\rceil)^2$ and the next squares: concretely
$(p-1)/2+1$ is not an integer square root situation --- $p+1$ is even and
$(\tfrac p2)^2<p+1<(p+2)^2$ with no square in between only for
$p\in\{3\}$; since $p\ge7$, $n$ is not a perfect square (if
$n=(m-1)(m+1)$ with $m\ge2$, then $p=(m-1)(m+1)$ is divisible by both
$m-1$ and $m+1$, so $m=2$ and $p=3$, contradiction). All three allowed
regimes die.
\end{proof}

\begin{remark}[scope of Theorem~\ref{thm:claprimefam}; binding]\label{rem:claprimscope}
Three boundary sentences travel with Theorem~\ref{thm:claprimefam}.
\emph{First}, it covers the non-infinity regimes only: the infinity
branch of the prime family is \emph{open} --- there is no paper proof that
infinity-mode cannot be sharp for $n=p+1$ with $p\equiv3\pmod4$, $p\ge7$
(the machine route through the layered arithmetic is blocked at the first
prime layer $d=p$ for $p\ge19$, where the fresh-enumeration volume
$\binom{3p-1}{2p-1}$ exceeds executability; $p=7$ and $p=11$ have been
completed, the latter closing the infinity branch of $K_{12}@66$). The
tower's infinity-mode deaths cover only $n=5..13$. Consequently the
extrapolations registered in the source material --- ``$\mathcal{E}$
contains the infinite family $\{K_{p+1}\}$'' and ``seed propagation
across $n=6..12$'' --- are \emph{withdrawn}: the theorem gives the
non-infinity subdomain only, is a sufficient-family statement and not a
classification, and no infinite-family lower bound for $\mathcal{E}$
follows from it until the infinity branch is closed. \emph{Second}, the
full per-$n$ verdicts require the infinity branch separately even where
the non-infinity side is dead ($n=8$ is closed by census; $n=12$ by the
completed layered enumeration; the remaining prime-family members are
open). \emph{Third}, the machine companion table for $n\le200$ lists $28$
full-death points, of which $23$ have $n-1$ prime (matching the theorem)
and $5$ are composite points $\{112,116,124,160,188\}$ --- the theorem is
a sufficient family, not a classification of the death set.
\end{remark}

\subsubsection{The corrected Fourier identity of the regime-general
character equation}\label{subsubsec:fourier}

The layered arithmetic used above has a closed-form face via the
characters of $\mathbb{Z}_k$. For a nontrivial character $\psi$ of
$\mathbb{Z}_k$, write, for a placement $S$ with finite points parametrized
as $y=2v$ (even) and $y=2w+1$ (odd),
\[
\Psi_E=\textstyle\sum_i \psi(v_i),\quad
\Psi_{E,2}=\textstyle\sum_i \psi(2v_i),\quad
\Psi_O=\textstyle\sum_a \psi(w_a),\quad
\Psi_{O,2}=\textstyle\sum_a \psi(2w_a).
\]

\begin{theorem}[corrected Fourier identity]\label{thm:clafourier}
Let $\psi$ be a nontrivial character of $\mathbb{Z}_k$.
\begin{enumerate}[label={\rm(\roman*)},leftmargin=2em]
\item (non-infinity placements) If a placement is sharp, then
\[
\tfrac12\bigl(\Psi_E^2-\Psi_{E,2}\bigr)
+\psi(1)\cdot\tfrac12\bigl(\Psi_O^2-\Psi_{O,2}\bigr)
+\Psi_E\Psi_O \;=\;0 .
\]
\item (infinity placements) If a placement with $\infty\in S$ is sharp,
then the infinity-class contribution is
\begin{equation}\label{eq:clarho}
\rho_\psi\;=\;\sum_i \psi(y_i \bmod k)\;=\;\Psi_{E,2}\;+\;\psi(1)\cdot\Psi_{O,2},
\end{equation}
i.e.\ both sides are doubled: the odd-side contribution is the doubled
half-value sum $\Psi_{O,2}=\sum_a\psi(2w_a)$, not $\Psi_O$. Equivalently,
the infinity-mode equation takes the ``$i\le j$'' form
\[
\tfrac12\bigl(\Psi_E^2+\Psi_{E,2}\bigr)
+\psi(1)\cdot\tfrac12\bigl(\Psi_O^2+\Psi_{O,2}\bigr)
+\Psi_E\Psi_O \;=\;0 .
\]
\end{enumerate}
\end{theorem}

\begin{proof}
(i) The classes of even--even pairs are
$\mathrm{class}(\{2v_i,2v_j\})=v_i+v_j$, of odd--odd pairs
$\mathrm{class}(\{2w_a+1,2w_b+1\})=w_a+w_b+1$, and of mixed pairs
$\mathrm{class}(\{2v_i,2w_a+1\})=v_i+w_a$ \eqref{eq:claclass}. Summing
$\psi$ over all unordered pairs with distinct classes, using the
algebraic identity
$\sum_{i\ne j}\psi(x_i+x_j)=\Psi^2-\Psi_{2}$ folded to $i<j$ (which
produces the $-\Psi_{E,2}$ resp.\ $-\Psi_{O,2}$ halves and the factor
$\psi(1)$ on the odd side from the $+1$ shift), the sharpness hypothesis
(class multiset $=\mathbb{Z}_k$ exactly) makes the total vanish.

(ii) The infinity edges $\{\infty,y\}$ have classes $y\bmod k$, so their
character contribution is $\sum_i\psi(y_i\bmod k)$. For an even point
$y=2v$ this is $\psi(2v)$; for an odd point $y=2w+1$ it is
$\psi(2w+1)=\psi(2w)\psi(1)$. Summing gives
$\rho_\psi=\Psi_{E,2}+\psi(1)\Psi_{O,2}$ \eqref{eq:clarho}. Substituting
into (i) and merging the doubled contributions --- each half-value sum now
appears both through its square and through its doubled self-sum --- turns
the ``$i<j$'' form into the ``$i\le j$'' form stated.
\end{proof}

\begin{remark}[why the corrected form is quoted; the falsified
predecessor]\label{rem:clafourierfix}
An earlier draft of the identity stated the infinity contribution as
$\rho_\psi=\Psi_{E,2}+\psi(1)\Psi_O$, omitting the odd-side double angle.
That form is false: on the $K_3$ configuration $\{\infty,1,3\}$ at $k=3$
with $\psi(a)=\omega^a$, $\omega^3=1$, $\omega\ne1$ (a sharp placement,
the classes being $1,0,2$), the original form gives the residual
$\omega+2\omega^2=-2-\omega\neq0$, while the corrected form gives $0$.
This was verified independently in exact $\mathbb{Z}[\omega]$ integer
arithmetic, and the corrected identity \eqref{eq:clarho} was verified
identically on all infinity placements at $k=3,6,10$ ($15+220+4{,}845$
placements; the original form fails on $8$, $178$ and $4{,}514$ of them
respectively). Two accompanying sentences bind every citation of the
layered-arithmetic machine results: \emph{all machine results (the tower
counting rules and the censuses) are independent of the character-equa
formula} --- the machine enumeration counts infinity classes by their
actual residues $2a$ and $2a+1$, which is the corrected bookkeeping ---
and the erroneous formula, wherever it existed, appeared only in the
internal source report and in no form in the frozen v2.5 submission.
\end{remark}

\subsubsection{The revival-depth law and the falsification of the
forall-\texorpdfstring{$k$}{k} target}\label{subsubsec:revdepth}

We now turn to the polar-corner family of the placeability programme. For
$k\ge5$ let $T_k=(k-1)\cdot K_2$ (the \emph{polar-corner family}): the
placement engine places the $k-1$ edges one by one, a \emph{state}
$\sigma=(U,C)$ at slot $j$ has $j$ components placed with occupied points
$U$ ($|U|=2j$) and occupied classes $C$ ($|C|=j$), free points
$F=\mathbb{Z}_{2k}\setminus U$, free classes $C_f=\mathbb{Z}_k\setminus C$,
and $\rho:=m-j=(k-1)-j$ unplaced components. The state is \emph{dead} if
$\mathrm{Span}(F)\cap C_f=\emptyset$ (no further $K_2$ can be placed); a
\emph{$d$-revival} of a dead state is the retraction of some $d$ placed
components (a set $W\subseteq U$ of points with $|W|=2d$ and a set
$Y\subseteq C$ of classes with $|Y|=d$) after which the remaining
$\rho+d$ components can all be placed legally; $d^*$ is the minimum $d$
admitting a revival.

\begin{theorem}[revival-depth lower bound]\label{thm:clarevdepth}
Let $T_k=(k-1)\cdot K_2$ and let $\sigma=(U,C)$ be a dead state at slot
$j$ with $\rho=(k-1)-j$ unplaced components. If $\sigma$ has a
$d$-revival, then $\rho+d\le3d$, i.e.\ $\rho\le2d$, i.e.
$d^*\ge\lceil\rho/2\rceil$.
\end{theorem}

\begin{proof}
Five steps, using only the placement semantics and the death identity.

(1) A $d$-revival retracts $d$ placed components; since every component of
$T_k$ is a $K_2$, the retracted point set $W$ has $|W|=2d$, and since each
placed component uses exactly one class, the retracted class set $Y$ has
$|Y|=d$.

(2) After the retraction, the $\rho+d$ components of the suffix must be
placed on the points $F\cup W$ using pairwise-distinct classes from
$C_f\cup Y$; the placements are pairwise point-disjoint with pairwise
distinct classes, i.e.\ they form a rainbow matching of size $\rho+d$ on
the board $(F\cup W,\,C_f\cup Y)$.

(3) Every edge $e$ of this matching either lies entirely in $F$ (an
``$F$-internal'' edge) or meets $W$; this dichotomy is exhaustive.

(4) An $F$-internal edge has its class in $\mathrm{Span}(F)$ (both
endpoints in $F$); but the death identity gives
$\mathrm{Span}(F)\cap C_f=\emptyset$, so its class lies in $Y$. Distinct
matching edges have distinct classes, so the number of $F$-internal edges
is at most $|Y|=d$.

(5) An edge meeting $W$ consumes at least one point of $W$, and matching
edges are disjoint, so their number is at most $|W|=2d$.

Steps (3)--(5) cover all edges of the matching, so
$\rho+d\le d+2d=3d$, i.e.\ $\rho\le2d$.
\end{proof}

\begin{remark}[the companion machine law and the paper-ization gap]
\label{rem:cladstarlaw}
The reverse inequality $d^*\le\lceil\rho/2\rceil$ is a \emph{machine law,
not a theorem}: it holds with zero violations over $4{,}273$ case slots
($4{,}252$ pool dead ends plus $21$ polar minima; deduplicated $4{,}268$),
but no paper proof of the upper-bound direction is known, and the
paper-ization of this law is registered as an open gap. The companion
deficit law $\min\text{-}\mathrm{deficit}(d)=\max(0,\rho-2d)$ is likewise
a machine law.

Consequence (falsification): the long-standing target statements $K_1$ and
the invariant refinement $[\mathrm{INV}\text{-}\mathrm{BT}''\text{-}2]$ in
$\forall k$ form are \emph{false}. The polar-corner minima at
$k=23,\dots,29$ have $\rho=5$, so by Theorem~\ref{thm:clarevdepth} no
$2$-revival exists; the machine channels agree (slot-pair channels $0$
out of $136$ and $0$ out of $153$ candidates, full $(W,Y)$-space sampling
$0$ out of $3{,}000$), and the former $60$-million-node ``double OPEN''
search walls at $k=23/24$ are structural, decided false in constant time
by the four-layer pruning of which Theorem~\ref{thm:clarevdepth} is the
last layer, with $d^*=3$ witnesses verified. The replacement targets are
the stratified statements: $d^*=\lceil\rho/2\rceil$ on the polar family
(lower bound theorem, upper bound machine law).
\end{remark}

\subsubsection{The criterion-level results: stratified statements only}
\label{subsubsec:criteria}

The remaining material of this section concerns criterion-level results:
statements about necessary conditions, their completeness on tested
domains, and their stratified theorem-form closures. A recurring discipline
binds all of them, which we state once here: \emph{a min-val$\le1$-type
criterion verified by per-instance machine scans over finitely many $N$ is
not an existence proof for general $N$}, and no sentence below may be read
as one.

\paragraph{The closed-form criterion family for the packing residual.}
For the $\mathbb{Z}_7$ packing problem at the top layer $\Sigma d=|S|$,
define the power-sum moment conditions $P_T$ ($\bmod\ 7$) and their
integer upgrades $\widetilde P_T$: there exist $C\subseteq S$ with
$|C|=\Sigma d$ and pairwise disjoint key sets $K_i\subseteq
\mathbb{Z}_7\setminus\{0\}$, $|K_i|=d_i$, matching all power sums up to
degree $T$. Necessity is immediate (a feasible packing makes the two sides
the same multiset), and at the top layer with $m=\Sigma d=|S|\le6$,
sufficiency of $P_m$ follows from Newton's identities ($t\le6$ invertible
in $\mathbb{F}_7$) forcing equality of characteristic polynomials, hence
of multisets: \emph{completeness at $T=|S|$} is a paper theorem with
machine backing ($3{,}872$ top-layer points, zero counterexamples). On the
census domain ($120$ orbits $\times\,330$ demand vectors $=39{,}600$
points, exhausted), the second integer moment $\widetilde P_2$ and the
third mod-$7$ moment $P_3$ agree with the oracle on every point ---
$0$ misclassifications out of $39{,}600$, zero false positives on all
$6{,}649$ feasible points --- and the residual core collapses $424\to0$ in
this domain (coverage ladder $1{,}133\to1{,}519\to1{,}557$). Out-of-domain
evidence is a sample, not an enumeration, quoted only at the corrected
scale: $1{,}200{,}000$ draws yield $33{,}018$ valid sample points
($31{,}242$ after deduplication; sampling support $933{,}450$;
$\approx3.5\%$ of support). Two boundary sentences bind any citation: the
$T=3$/$T=2$ closure is a domain-exhaustion statement (the census domain),
not a theorem --- the only paper-level completeness is at $T=|S|\le6$ ---
and the integer moments are not affinely equivariant, so cross-domain
citations use the mod-$7$ version. A convex-unreachability companion
holds in its corrected scope: if an infeasible point is a rational convex
combination of feasible colour-count vectors ($38/38$ explicit
certificates, exactly re-checked), then no criterion that is a relaxation
of the natural fractional LP can exclude it --- its original
``any-rank Chv\'atal--Gomory'' listing clause is refuted by rank-$1$
CG single cuts ($9/38$ of the very certificates) and is never quoted.

\paragraph{The stratified valuation criterion.}
The valuation criterion in its aligned family reads: \emph{for every
instance with the three anchor offsets on a common coset of the
subgroup generated by the first anchor difference, and with capacity slack
$\beta\ge2$, the minimum of the valuation $\mathrm{val}(C_1,C_2)$ over
admissible pairs is $\le1$}; the machine scan domain is $\beta\ge1$. We
register the criterion and its existence half as formal objects:
\begin{conjecture}[valuation criterion, candidate form]\label{conj:clacand}
Every so1 instance with $\beta\ge2$ has $\min\mathrm{val}\le1$.
\end{conjecture}
The statement is a \emph{candidate}: no paper proof exists, its
falsification clause is an explicit so1 instance with $\beta\ge2$ and
$\min\mathrm{val}>1$, and every machine account of this line is a directed
sweep of exactly that clause, with zero hits on the tested domains. The
general-$N$ existence half is carried separately:
\begin{conjecture}[the residual existence gap]\label{int:conj-gapex}
For every $N$ and every so1 instance with $\beta\ge2$,
$\min\mathrm{val}\le1$.
\end{conjecture}
The status of Conjecture~\ref{conj:clacand} is stratified, and only this
stratified form may be cited:
\begin{itemize}[leftmargin=1.5em]
\item \emph{Per-$N$ closure in theorem form} for $N\in\{6,8,10,12\}$: a
three-tier finite search, proved complete and terminating on paper (a
witness-bound theorem gives the per-instance bound $B(I)=d_2+1$ and a
uniform $B(N)=N-2$; a completeness theorem shows the three searched
families cover every witness), with zero residue on
$588/30{,}760/634{,}960/7{,}798{,}903$ instances respectively.
\item \emph{The general-$N$ existence half of
Conjecture~\ref{int:conj-gapex} is open}: whether for all $N$
and all instances with $\beta\ge2$ an admissible pair with
$\mathrm{val}\le1$ exists. The $C_1$ side is closed for all $N$ by a
constructive one-leak theorem --- for every instance and every admissible
$d_1$ there is $C_1$ with at most one leak, by an explicit three-case
construction over the cycle decomposition (full-cycles, partial cycles,
and the cycle containing the first anchor) --- so the obstruction is
entirely on the $C_2$ side (reachability of a near-closed board under the
capacity identity). The phrase ``theorem-level closure of the criterion''
is an explicit non-claim, including at larger $N$.
\end{itemize}

\paragraph{The covering identity as a candidate.}
The strongest structural statement of the valuation line is the following
six-family identity, quoted strictly at candidate tier.

\begin{proposition}[covering identity; paper candidate assembled and
independently re-audited; not mechanized]\label{prop:clacover}
For every instance $I=(N,P,T,d)$ with $P\subseteq\mathbb{Z}_N$,
$|P|\ge4$, $T=(t_1,t_2,t_3)$ pairwise distinct,
$d\in\mathbb{Z}_{\ge0}^3$, $\Sigma d_i=|P|-1$,
$d_i\le|P|-[t_i\in P]$: let $W=P\setminus\{t_1\}$,
$k=d_2+d_3+[t_1\notin P]$ (so $|W|-k=d_1$),
$\delta_{12}=(t_1-t_2)\bmod N$, $\delta_{13}=(t_1-t_3)\bmod N$, and
$L=\{R\subseteq W:\ |R|=k\}$. Let $\mathrm{Im}$ be the set of position-sets
produced by the eight-branch family generator, and let
$C_{12},C_{13},\mathrm{RX},O3_{12},O3_{13},V_7$ be the six covering
families defined via the per-orbit run statistic $\mathrm{rc}_\delta$
(full containment counts $1$), cross-reachability, full-orbit removal, and
the image of the legacy candidates routine. Then
\[
\mathrm{Im}\cap L \;=\; (V_7\cap L)\cup(C_{12}\cap L)\cup(C_{13}\cap L)
\cup(\mathrm{RX}\cap L)\cup(O3_{12}\cap L)\cup(O3_{13}\cap L).
\]
\emph{Corollary (boundary).} $R\in L\setminus\mathrm{Im}$ implies
$\mathrm{rc}_{12}(R)\ge3$, $\mathrm{rc}_{13}(R)\ge3$, not
cross-reachable, not $O3_{12}$, not $O3_{13}$, and $R\notin V_7$.
\end{proposition}

\begin{remark}[tier and riders of Proposition~\ref{prop:clacover}]
\label{rem:clacoverriders}
The identity is a paper candidate (both directions have paper proofs and
the package passed an independent from-scratch re-derivation audit); it is
\emph{not} a mechanized or formalized theorem, and \emph{no disjointness
of the six families is claimed}. The machine-checkable face is the
five-family identity without the $V_7$ term: all machine checks are run in
the \emph{no-$V_7$ reading} (the $V_7$ branch omitted, giving a
conservative superset $\mathrm{fam}_{\mathrm{no}V_7}\subseteq
\mathrm{Im}\cap L$, a one-way implication toward the true face), and the
$V_7$ branch's machine face is frozen --- the legacy candidate-generation
source files were lost with no recoverable copy --- while the $V_7$ term
itself is definitional. The residual $V_7$-only mass is nonzero and
precisely measured: $340$ positions at $N=8$ ($1{,}266$ superset-uncovered
minus $926$ true-face uncovered, both sides measured) and $2{,}267{,}655$
instance-weighted at $N=14$ ($52{,}572{,}256$ minus $50{,}304{,}601$).
The supporting lemma list ($L1$--$L9$) is paper-proved; $L1/L9$ are
definitional ($V_7$ outputs), and the count-difference identity inside
$L8$ counts whole \emph{orbits}, not points (a point-count reading
mismatches in $136/936$ enumerated configurations). One reverse clause of
the run-statistic calculus must be stated one-way: ``$\mathrm{rc}=2$
implies a union of two disjoint cyclic pieces; the converse only gives
$\mathrm{rc}\le2$''. The three micro-lemmas (single interval
$\Rightarrow\mathrm{rc}=1$; two disjoint intervals
$\Rightarrow\mathrm{rc}\le2$; subadditivity) carry paper proofs plus two
rounds of randomized machine checks with zero violations ($12{,}558 /
2{,}500 / 12{,}000$ first pass; $4{,}000/1{,}109/4{,}000$ audit pass).
\end{remark}

\paragraph{The falsification-scale facts of the exchange ledger.}
The single-swap ledger statement --- that an uncovered position admits an
in-family companion at swap distance $2$ --- does \emph{not} hold
universally. On the fixed random pool ($900$ instances, $300$ each at
$N\in\{10,12,14\}$; $1{,}501$ uncovered-witness positions in the no-$V_7$
superset reading): a distance-$2$ companion in the family image exists for
$1{,}469/1{,}501$ ($97.87\%$) of witnesses, and within the five-family
domain for $1{,}447/1{,}501$ ($96.40\%$). Failing witnesses: $54/1{,}501$
($3.60\%$), concentrated in $12/900$ instances ($1.33\%$); all $12$ lie in
$N\in\{12,14\}$ with $|P|$-spectrum
$\{5,6,9,9,10,11,11,12,12,12,12,14\}$ --- $10$ of $12$ biased to high
$|P|$, $2$ outside the dense band. The cause split of the $54$ is
$32$ structural-side (no in-family companion exists) against $22$
ledger-side (companions exist but none satisfies the ledger constraints).
The distance-$\le4$ binary form also fails on the same pool: $26/54$ have
a distance-$4$ five-family companion, $28/54$ do not. Two accounting
disciplines bind every quotation: the instance-level ($12/900$) and
witness-level ($54/1{,}501$) figures are different measurement layers and
are never interchanged, and the $96.40\%$ figure is a
no-$V_7$-superset-reading statistic not directly comparable to
true-reading protocols.

\begin{remark}[the bridge guards]\label{rem:clabridge}
The bridge between the $\mu$-sufficiency question (whether satisfaction of
all $\mu$-inequalities implies feasibility on the difference board) and
the valuation criterion carries three explicit guards, all of which must
accompany any quotation. \emph{(a)} The equivalence ``abstract theorem
$\iff$ $\mu$-sufficiency on the difference board'' holds only under the
$(H2)/n\le p-1$ guard; the unguarded form is falsified by an explicit
$\mathbb{Z}_{14}$ counterexample ($P=\{0,1,2\}$, $T=(0,1)$, $d=(1,2)$:
all $\mu$-inequalities hold, the board is infeasible). \emph{(b)} The
identification of the abstract existence gap with the general-$N$
existence half above is valid only as a $k=3$-layer monotone-reduction
pointer; as a literal identity it fails on four axes (the $\beta$-band,
the so-slice, the $k$-domain, and the closure grade). \emph{(c)} The
$k\ge4$ layer is not touched by the criterion at all --- no criterion
exists in the current language for $k\ge4$, a language-level gap rather
than an unexecuted proof. Consequently: a min-val$\le1$-type criterion
verified by per-instance machine scans over finitely many $N$ is not an
existence proof for general $N$, and $\mu$-inequality checks over
finitely many $N$ do not constitute an existence theorem. Necessity of
the $\mu$-inequalities is strictly proved (feasible placements satisfy
all of them), the sufficiency direction is false on general bipartite
boards (an explicit $2\times2$ counterexample whose colouring violates
the difference rule), and sufficiency on the difference board is open
(zero violations in $N=6..14$ machine accounts; no general-$N$ paper
proof). Finally, the $\mu$ $k=2$ stratified existence question --- the
main open wall of this line --- is open in its paper face (machine face
closed on seven domains for the $k=2$ layer; recorded obstruction: no
base-to-step mechanism), and it has \emph{no logical dependency} on the
family-image characterization: characterizing the family image is not the
realize-layer existence.
\end{remark}

\begin{remark}[what the criterion layer does and does not give]
\label{rem:clacriterionlimits}
The criterion-level results of this subsection share one logical shape:
each is a necessary condition, or a family of necessary conditions, whose
completeness is established on explicitly enumerated domains, or whose
per-$N$ verification is theorem-form on the listed $N$. None of them
yields, unconditionally and for general parameters, the existence of the
object the criterion is meant to witness: the valuation criterion does
not give feasible pairs for general $N$ (the existence half is open), the
moment criteria do not give packings outside the census domain (and the
top-layer completeness is bounded at $|S|\le6$), and the covering identity
characterizes an image, not a realizability theorem. In particular the
conclusion of the conditional chain consuming these criteria remains
conditional: no unconditional existence statement, and hence no
unconditional progress on the underlying decomposition conjecture, is
claimed anywhere in this section.
\end{remark}
